% 【理论扩展】所属：第5章（由 chapters/revised/05-social-knowledge.tex \input 引入）
\minitopic{社会认识论的两种取向}一种取向以个人为最终主体，研究证言、分工和网络怎样影响个人真信念；另一种承认群体、制度具有不可简单还原的信念、知识和责任。前者可称个人主义社会认识论，后者研究集体认识主体。两者共享问题，却选择不同分析单位。 Goldman以真理促进评价社会实践\parencite{goldman1999}。批评者担心，只数真信念会忽略理解、权力和公正。更全面评价可包含准确性、纠错、信息覆盖、参与和损害分配。指标多元会产生权衡，但比假定单一目标更透明。 社会性也不表示相对主义。某共同体一致相信$p$不能使$p$为真；社会机制仍可按是否发现真相评价。研究社会条件恰恰解释客观知识如何在有限、偏见个体之间可能。

\minitopic{认识分工与依赖}现代科学、法律和新闻依赖专业分工。没有人亲自验证供应链、仪器、统计与全部文献。Hardwig指出，理性有时要求相信自己无法完全理解的专家证言\parencite{hardwig1985}。认识自主因而不能等于独立重做一切。 分工的收益是深度和效率，风险是接口失败。实验人员知道数据质量，统计人员知道模型，决策者可能只见摘要。若每层压缩不确定性，最终结论会显得比基础证据更确定。良好接口保留来源、假设和异常。 冗余可提高安全，但完全重复浪费资源。异质冗余更有价值：不同方法、数据和团队独立汇合，能发现共同盲点较少。制度应区分真正复制与同一流程的重复运行。

\minitopic{群体信念的形成}多数投票不是群体信念的唯一模型。董事会可按章程正式接受一项计划，即使部分成员私下反对；研究团队可在报告中承诺结论，形成对外可问责态度。群体信念与成员信念相关却不等同。 功能观点要求群体能以结论进行推理、行动、响应反证。档案中孤立文件不使组织知道，除非信息进入可调用流程。反之，没有任何成员掌握全部内容的分布式系统，也可能通过角色协调实现这些功能。 群体知识还需真实性与适当形成机制。组织正式相信真命题，却由篡改流程碰巧得真，是群体Gettier案例。成员善意不修复制度运气；需要审计和纠错机制。

\minitopic{判断聚合与话语悖论}三位委员分别判断命题$p$、$p\rightarrow q$和$q$。多数支持$p$，多数支持$p\rightarrow q$，却多数反对$q$，群体判断不逻辑闭合。这是话语困境：逐命题多数与按个人整体立场投票可能得不同结果。 制度必须选择聚合对象。逐议题投票尊重每个判断的多数，却可能不一致；先让各方提出完整方案再投票，保持一致却可能在单项上违背多数。没有规则同时满足所有吸引人的条件。 公开论证可在投票前改变个人判断，减少纯聚合困难，但权力和发言差异会影响讨论。形式结果不是说民主不理性，而是提醒规则设计有实质认识后果。

\minitopic{科学共同体与批评结构}Longino认为，科学客观性来自转化性批评：有公开批评渠道、共享标准、共同体响应以及一定程度的智识权威平等\parencite{longino1990}。个体价值背景不可完全消除，但多元批评能揭示隐藏假设。 同行评议只是其中一环。匿名可减轻地位偏见，也降低问责；开放评议增加透明，却可能压制年轻研究者批评权威。制度选择应依据领域规模、报复风险和可验证性，不存在统一最优。 复现失败不自动证明欺诈。样本差异、方法模糊、统计波动和选择性报告都可能。共同体要记录失败、分析异质性并修正理论，而不是把复现当简单有罪判决。

\minitopic{立场与情境知识}立场认识论指出，社会位置影响主体接触何种事实和概念。照护者、基层工人或被监管群体可能看到决策者忽略的机制。位置提供认识机会，却不保证自动正确；经验仍需解释、比较和批评。 “局内人最懂”与“专家最懂”都过于绝对。局内人有体验与情境接触，专业研究者有比较工具和方法训练。合作可结合优势，前提是局内证言不被只当原始数据，专家也不放弃方法责任。 情境性不等于真理只对群体成立。它说明证据可得性和问题选择受位置影响。扩大参与能发现遗漏事实，同时结论仍可接受跨位置检验。

\minitopic{认识不正义的扩展}证言不正义聚焦可信度亏欠，但身份偏见有时也给予过度可信度，如默认某类人天然权威。过度信用可能让主体承担不愿代表群体的负担，也使错误不受审查。公正要求信用与相关能力和证据相称。 诠释资源缺口不仅缺一个词，还可能缺少让经验进入制度的分类、表格和程序。新的概念能使经验可交流，却也可能被行政系统僵化。受影响者应参与概念形成与修订。 Dotson讨论认识暴力与沉默：说话者有内容和能力，但听者因无能或拒绝而使交流失败\parencite{dotson2011}。解决不只训练说话者更清楚，还需改变听者习惯和权力结构。

\minitopic{回音室、气泡与极化}认识气泡因网络结构遗漏某些来源；加入缺失声音可能修复。回音室更主动地贬低外部来源，成员把反证解释为敌对证明。后者不能只靠增加信息量解决，需要重建跨边界信任。 群体极化可能来自论据池不平衡、社会比较和身份信号。讨论成员多不等于观点多；若都依赖相同来源，交流会重复强化。设计可要求独立先行判断、轮流反方和来源谱系展示。 平台推荐常优化参与而非真理。高情绪内容获得更多反馈，改变可见证据。个人媒介素养重要，但平台目标、转发摩擦和来源标识同样是认识制度问题。

\minitopic{民主的认识价值}民主可因汇集分散信息、允许公开批评和和平纠错而有认识优势；多数也会受宣传、无知和协调失败影响。认识论辩护不应假定多数永远正确，而应比较制度发现和修正错误的能力。 专家机构与公众参与不是简单对立。技术问题需专业判断，政策目标与风险分配需民主正当性。混合制度可让专家公开证据和不确定性，受影响群体提出遗漏问题，决策者说明价值权衡。 透明本身也有边界。公开所有讨论可能使成员表演立场、阻碍试探性意见；完全保密则隐藏利益。分阶段透明、记录理由和事后审计可在学习空间与问责之间平衡。

\minitopic{综合案例：城市空气质量报告}监测站由承包商运行，大学团队建模，政府发布指数，居民报告气味与症状。某日官方指数正常，社区却普遍不适。不能预先断定仪器客观、居民主观，也不能因居民众多就否定模型。 调查需检查站点分布、污染物类别、时间平均是否掩盖峰值、居民报告是否独立、模型对该社区是否校准。社区位置可能揭示未测暴露，专业方法能排除其他原因。共同设计新增监测可把冲突转成可检验问题。 若历史上监测站长期远离低收入社区，居民不信任有记录基础。重建信任需要改变采样与治理，让数据、维护记录和申诉可见。最终知识属于一个经修正的群体过程，而非某方先天特权。

\begin{tcolorbox}[breakable,colback=white,colframe=blue!30!black,title={社会认识系统审计表},fonttitle=\bfseries]
列出主体、角色和信息接口；追踪来源独立性与压缩损失；检查群体态度形成规则和逻辑一致性；评价批评、复现与纠错渠道；分析可信度分配和诠释资源；识别激励、网络和权力；同时报告真理表现、公正与可问责性。
\end{tcolorbox}
